\textbf{Example 1}

Consider the following call:

\texttt{construct\_roads([4, 4, 6, 4, 2], [4, 6, 4, 2, 4])}

This means that there are $5$ fountains:
\begin{itemize}
\item fountain $0$ is located at $(4,4)$,
\item fountain $1$ is located at $(4,6)$,
\item fountain $2$ is located at $(6,4)$,
\item fountain $3$ is located at $(4,2)$,
\item fountain $4$ is located at $(2,4)$.
\end{itemize}

It is possible to construct the following $4$ roads, where each road connects two fountains, and place
the corresponding benches:

\begin{center}
\renewcommand{\arraystretch}{1.5}
\begin{tabular}{|c|c|c|}
\hline
Road label & Labels of the fountains the road connects & Location of the assigned bench \\
\hline
$0$&$0,2$&$(5,5)$\\
\hline
$1$&$0,1$& $(3,5)$\\
\hline
$2$&$3,0$& $(5,3)$\\
\hline
$3$&$4,0$& $(3,3)$\\
\hline
\end{tabular}
\end{center}

This solution corresponds to the following diagram:

\includegraphics[scale=0.7]{1.png}

To report this solution, \texttt{construct\_roads} should make the following call:
\begin{itemize}
\item \texttt{build([0, 0, 3, 4], [2, 1, 0, 0], [5, 3, 5, 3], [5, 5, 3, 3])}
\end{itemize}

It should then return $1$.

Note that in this case, there are multiple solutions that satisfy the requirements, all of which would be
considered correct. For example, it is also correct to call \texttt{build([1, 2, 3, 4], [0, 0, 0, 0], [5, 5, 3, 3], [5, 3, 3, 5])} and then return $1$.

\textbf{Example 2}

Consider the following call:

\texttt{construct\_roads([2, 4], [2, 6])}

Fountain $0$ is located at $(2,2)$ and fountain $1$ is located at $(4,6)$. Since there is no way to
construct roads that satisfy the requirements, \texttt{construct\_roads} should return $0$ without making any call to \texttt{build}.
